\section{Task 4: The test lines and the accuracy on real photographs}
\sectionmark{Task 4: Test lines and accuracy}
\label{sec:t4}

\subsection{Purpose, conditions and criterion}

\begin{plain}{Task 4 at a glance}
\textbf{Purpose.} The accuracy figures known so far come from simulated photographs. This task
replaces them by measurements on real ones. In part~A the 80 printed lines are photographed
with the board. In part~B, after the pipeline has been frozen, the 60 test lines among them are
evaluated once, and the result is compared with the targets of the project.

\textbf{Before you begin.} For part~A: Task~1 is complete, the line platform exists, and the 16
printed pages of test lines are at hand. For part~B: Task~3 is complete, and the frozen version
carries the tag \texttt{v1.0-rc1}.

\textbf{Complete when.} Part~A: 80 photographs exist, each shows the whole of the line after
which it is named, and they are stored in the repository. Part~B: the table of results against
targets (Table~\ref{tab:targets}) is filled in with counts, and the per-line results and the
review document are stored in the team's shared folder.

\textbf{Time.} Part~A about one hour; part~B about half an hour.
\end{plain}

\subsection{Part A: photograph the 80 lines}
\label{sec:t4a}

Part~A may be carried out while the classifier of Task~2 trains. In that case it needs a second
PowerShell window, and the window in which the training runs must be left alone.

\step{Open a PowerShell window for the photography}
Click \ui{Start}, type \code|powershell|, and click \ui{Windows PowerShell}. A new window opens.
In it, give the two commands of the opening sequence:
\begin{ps}[Opening sequence, command 1 of 2]
cd C:\ee105\camera-to-latex
\end{ps}
\begin{ps}[Opening sequence, command 2 of 2]
.\.venv\Scripts\Activate.ps1
\end{ps}
\checkpoint{The prompt of the new window reads:}
\venvprompt

\step{Set up the line position}
Put the line platform back on its outline, with a blank white sheet on top, and check the lamps
as in Task~1. Sort the 16 printed pages by the names at their right margins: the pages with T01
to T60 first, then those with D01 to D20.
\checkpoint{The platform has the thickness entered in the record of Task~1, and the pages are
in order.}

\step{Start the series}\label{st:lineseries}
\begin{pswide}
python -m pipeline.capture --series test_set\images --names-file test_set\photo_names.txt
\end{pswide}
The viewfinder opens, and its yellow line reads \emph{1/80\ \ next: T01}. The names follow the
order of the printed pages.
\checkpoint{The viewfinder shows the live picture and the name \emph{T01}.}

\step{Photograph each line}\label{st:oneline}
For each of the 80 lines:
\begin{enumerate}
\item Lay the page on the platform, its top towards the arrow, and move it until the line named
  in the viewfinder lies in the middle of the picture. The name is printed at the right margin
  of the page, level with the line, and is itself outside the picture.
\item Check four things in the viewfinder. The whole line is inside the picture, with white
  paper beyond its first and its last symbol. The line lies in the middle third of the height
  of the picture and runs parallel to the border, its two ends at the same height to within
  half a letter. No shadow and no bright spot touches it. For a fraction or a sum, the whole
  structure is inside, top to bottom.
\item Wait until the picture has stood still for two seconds, and press \key{Space}.
\end{enumerate}
Parts of the neighbouring lines appear at the upper and lower borders of the picture. They are
left as they are: the program reads the row nearest the middle, and covering the neighbours
would only add the edges and shadows of the covering sheets. Line T39 is the only one that
nearly fills the width of the picture; centre it with particular care.
\checkpoint{After the 80th photograph the viewfinder closes by itself.}

\begin{note}[Mistakes and interruptions]
As in Task~2: \key{B} goes back one name so that a photograph can be replaced; \key{Q}
interrupts the series; when the command is given again, \key{N} passes over the names marked
\emph{(exists)}. Single photographs are taken again with the placeholder form of the command,
in which the names are separated by commas and enclosed in quotation marks:
\begin{center}
\mbox{\footnotesize\texttt{python -m pipeline.capture -{}-series test\_set\textbackslash{}images -{}-names "}\ph{NAMES}\texttt{"}}
\end{center}
For example, \code|--names "T17,T23"| at the end of the command retakes the lines T17 and T23.
\end{note}

\step{Count and inspect the photographs}
\begin{ps}
(Get-ChildItem test_set\images -Filter *.png).Count
\end{ps}
\checkpoint{The answer is \texttt{80}.}
\begin{ps}
explorer test_set\images
\end{ps}
In the Explorer window choose \menu{View}\then\menu{Extra large icons} and look at all 80
photographs. Each must show one whole line in the middle of the picture, upright and sharp.
Then compare five photographs with the printed lines of the same names: T05, T25, T45, T60 and
D20.
\checkpoint{Every photograph shows a whole line, and the five photographs show the lines after
which they are named. A photograph that shows the neighbour of its line means that a line was
skipped or taken twice; from that point on the names are shifted, and the affected photographs
are taken again.}

\step{Let the pipeline read the development lines}
The 20 development lines, and only they, may be evaluated at once; the test lines wait for
part~B.
\begin{ps}
python -m tools.eval_pipeline --split dev --no-compile --quiet
\end{ps}
\checkpoint{The table reports 20 lines in its row \emph{all}, and a share of exact readings
that is clearly above zero. (The column \emph{compiles} shows \emph{nan}, because that check
was omitted.) A share of zero means that the photographs are unusable as a set:
the names are shifted, or the orientation setting of Task~1 is missing on this laptop.}

\step{Store the photographs in the repository}
\begin{ps}
git add test_set\images
\end{ps}
\begin{ps}
git commit -m "Photographs of the 80 test and development lines"
\end{ps}
\begin{ps}
git push
\end{ps}
\checkpoint{The output of the last command ends with a line that contains \texttt{main ->
main}.}

\step{Fill in the record of part A}
\begin{table}[H]
\centering
\begin{tabular}{L{7.4cm}C{3.8cm}}
\toprule
\multicolumn{1}{c}{\hd{Quantity}} & \hd{Value} \\
\midrule
Date; photographer & \\
Thickness of the platform; width $W$ of T39 in \file{T39.png} & \\
Lamps (number, position) & \\
Photographs taken; photographs retaken & \\
Print size of the lines & \\
\bottomrule
\end{tabular}
\caption{Record of Task 4, part A. The width of T39 is measured in Paint as in Task~1, with
the command \texttt{mspaint test\_set\bs images\bs T39.png}.}
\label{tab:t4arecord}
\end{table}

\subsection{Part B: the final measurement}
\label{sec:t4b}

The final measurement is made once, on the frozen version. Its result is reported whatever it
is; a target that is missed is reported as missed.

\step{Open the project in PowerShell}
% The opening sequence, written out in full. Input wherever a task may begin in a new window.
If a PowerShell window whose prompt begins with \texttt{(.venv)} is still open, use it and go
on to the next step. Otherwise open a new window (click \ui{Start}, type \code|powershell|,
click \ui{Windows PowerShell}) and give the two commands of the opening sequence in it.
\begin{ps}[Opening sequence, command 1 of 2]
cd C:\ee105\camera-to-latex
\end{ps}
\begin{ps}[Opening sequence, command 2 of 2]
.\.venv\Scripts\Activate.ps1
\end{ps}
\checkpoint{The prompt reads:}
\venvprompt


\step{Confirm that the frozen version is in use}
\begin{ps}
git pull
\end{ps}
\begin{ps}
git status --short
\end{ps}
\begin{ps}
git describe --tags
\end{ps}
\checkpoint{The second command prints nothing: no file differs from the repository. The third
prints \texttt{v1.0-rc1}. If it prints a longer name, such as \texttt{v1.0-rc1-2-g3f2c1ab},
commits have been added since the freeze; ask the lead whether they were intended.}

\step{Evaluate the 60 test lines}\label{st:final}
\begin{ps}
python -m tools.eval_pipeline
\end{ps}
The command reads the 60 photographs T01 to T60, compares each reading with the expected
\LaTeX{}, checks that each reading compiles, prints one line for each photograph and the table
of Step~\ref{st:gaterun}, and names the file in which the per-line results are stored.
\checkpoint{The row \emph{all} reports 60 lines. The table is identical to that of the gate of
Task~3, because neither the photographs nor the pipeline have changed and the pipeline contains
no element of chance. A difference means that something has changed since the gate; find out
what before any figure is reported.}

\step{Produce the review document}
\begin{ps}
python -m tools.review_results
\end{ps}
\begin{out}
review: C:\ee105\camera-to-latex\results\review_test_20261130_1012.pdf
\end{out}
\checkpoint{The file exists. It shows all 60 lines, each with its photograph, the boxes, the
expected line and the reading; misread lines are marked in orange.}

\step{Compare the result with the targets}
Convert each percentage of the table into a count, and fill in Table~\ref{tab:targets}. The
accuracy on single symbols comes from the record of Task~2, and the time per picture from the
soak test of Task~1.

\begin{table}[H]
\centering
\begin{tabular}{L{4.5cm}C{2.5cm}C{2.0cm}C{1.5cm}}
\toprule
\multicolumn{1}{c}{\hd{Quantity}} & \hd{Target} & \hd{Result} & \hd{Met?} \\
\midrule
Symbols named correctly (held-out real crops) & 95 \% & \hfill of 280 & \\
Plain lines read exactly & 80 \% (20 of 25) & \hfill of 25 & \\
Script lines read exactly & 70 \% (14 of 20) & \hfill of 20 & \\
Fraction lines read exactly & 60 \% (6 of 10) & \hfill of 10 & \\
Mixed lines read exactly & no target & \hfill of 5 & -- \\
All test lines read exactly & no target & \hfill of 60 & -- \\
Readings that compile & 100 \% & \hfill of 60 & \\
Time from key press to reading & 5 s or less & \hfill s & \\
\bottomrule
\end{tabular}
\caption{Results against the targets of the project plan. The time from key press to reading is
the time per picture of the soak test plus the value in the column \emph{s/line} of the
evaluation. In the application a line can be accepted only if it compiles, so that every
accepted line compiles by construction; the row on compiling reports how many readings
compiled exactly as they were produced.}
\label{tab:targets}
\end{table}
\checkpoint{Every row has a result. Each result is a count, so that its precision can be
judged (Section~\ref{sec:stat}).}

\step{Classify the misread lines}
Open the review document (\code|explorer results|, then double-click the file whose name begins
with \file{review\_test}) and assign every misread line to a kind of error by
Table~\ref{tab:kinds}, as in Step~\ref{st:triage}. Mark three or four lines that show different
kinds of error clearly; the final report presents them as examples.
\checkpoint{The record states how many misread lines fall under each code. This analysis
serves the report alone: nothing in the pipeline is changed because of it.}

\step{Store the evidence}
In the Explorer window of the folder \file{results}, sort by date. Upload the two newest files
whose names begin with \file{eval\_test} and \file{review\_test} to the team's shared Google
Drive folder (\menu{New}\then\menu{File upload}).
\checkpoint{Both files are visible in the shared folder. The final report takes its figures and
its examples from them.}

\step{Enter the result in the table of the README (project lead)}
The file \file{README.md} of the repository contains a table headed \emph{Results}, with one
row for each measurement.
\begin{ps}
notepad README.md
\end{ps}
Press \key{Ctrl}+\key{F}, type \code|## Results| and press \key{Enter}; Notepad jumps to the
heading of the table. Select the last row of the table, which is one long line that begins and
ends with a vertical bar, copy it, and paste the copy on a new line below it. In the copy,
replace the entries between the bars by today's date, the version \code|v1.0-rc1|, the words
\emph{60 real test lines}, the capital height of Task~1, the classifier, and the counts of
Table~\ref{tab:targets}. Keep every vertical bar. Save with \key{Ctrl}+\key{S}, close Notepad,
and record the change:
\begin{ps}
git add README.md
\end{ps}
\begin{ps}
git commit -m "Results of the final measurement on real photographs"
\end{ps}
\begin{ps}
git push
\end{ps}
\checkpoint{On the GitHub page of the repository, the table shows the new row.}

\step{Fill in the record of part B}
\begin{table}[H]
\centering
\begin{tabular}{L{7.4cm}C{3.8cm}}
\toprule
\multicolumn{1}{c}{\hd{Quantity}} & \hd{Value} \\
\midrule
Date; version (output of \texttt{git describe -{}-tags}) & \\
Name of the per-line results file & \\
Table identical to that of the gate? & \cbox\ yes\quad\cbox\ no \\
Misread test lines by kind: P; C1; C2; C3; R; L & \\
Lines chosen as examples for the report & \\
\bottomrule
\end{tabular}
\caption{Record of Task 4, part B. The results themselves are in Table~\ref{tab:targets}.}
\label{tab:t4brecord}
\end{table}
